\section{Autonomous Coding Agents}
\label{sec:agents}

Autonomous coding agents represent the frontier of AI-assisted development. Unlike assistants that require manual application of suggestions, these systems can directly read files, write code, execute commands, and iterate on their work.

\subsection{Defining Autonomy}

An autonomous coding agent has:
\begin{itemize}
    \item \textbf{Environment access}: Can read and write files in an actual codebase
    \item \textbf{Execution capability}: Can run commands, tests, and scripts
    \item \textbf{Feedback loops}: Observes results and iterates on errors
    \item \textbf{Multi-step planning}: Decomposes complex tasks into actionable steps
\end{itemize}

The key distinction from assistants is operational: agents \emph{act} on code rather than merely \emph{advising} about it.

\subsection{Architectural Patterns}

Most autonomous agents follow variations of an observe-think-act loop:

\begin{enumerate}
    \item \textbf{Observe}: Gather information (read files, search code, check test results)
    \item \textbf{Think}: Reason about the task and plan next steps
    \item \textbf{Act}: Execute an action (edit file, run command)
    \item \textbf{Repeat}: Process results and continue
\end{enumerate}

\textbf{Tool use} is central to agent capabilities. Common tools include:
\begin{itemize}
    \item File operations (read, write, search, navigate)
    \item Shell command execution
    \item Web search and documentation lookup
    \item Browser automation for testing
\end{itemize}

\textbf{Memory} mechanisms help agents maintain context across long tasks:
\begin{itemize}
    \item Conversation history
    \item Scratchpads for working notes
    \item Retrieved context from codebase search
\end{itemize}

\subsection{Major Systems}

\subsubsection{SWE-agent}

SWE-agent~\citep{yang2024sweagent} from Princeton introduces the Agent-Computer Interface (ACI)---a set of custom commands designed for LLM interaction with codebases. Rather than raw shell access, SWE-agent provides commands like \texttt{find\_file}, \texttt{search\_dir}, and \texttt{edit} with structured output formats optimized for LLM parsing.

Key innovations:
\begin{itemize}
    \item ACI reduces errors from shell command complexity
    \item Search/replace editing prevents whole-file regeneration
    \item Achieves 12.5\% on SWE-bench (original) with GPT-4
\end{itemize}

\subsubsection{Devin}

Devin~\citep{cognition2024devin} from Cognition Labs is positioned as an ``AI software engineer'' with access to a complete development environment: shell, browser, and editor. It can execute long-running tasks autonomously, maintaining state across sessions.

Key features:
\begin{itemize}
    \item Full development environment (not just file access)
    \item Planning and checkpoint system for complex tasks
    \item Browser for testing and documentation lookup
    \item Reported strong SWE-bench performance (proprietary evaluation)
\end{itemize}

\subsubsection{OpenHands}

OpenHands~\citep{wang2024opendevin} (formerly OpenDevin) provides an open-source platform for building coding agents. It offers:
\begin{itemize}
    \item Modular architecture supporting different agent implementations
    \item Sandboxed execution environment
    \item Web interface and API access
    \item Multiple agent implementations for research comparison
\end{itemize}

\subsubsection{Claude Code}

Claude Code~\citep{anthropic2024claudecode} is Anthropic's CLI-based coding agent. It emphasizes:
\begin{itemize}
    \item Terminal-native workflow for existing developer tools
    \item Permission system for controlled autonomy
    \item Integration with git workflows
    \item Extended thinking for complex reasoning
\end{itemize}

\subsubsection{Aider}

Aider~\citep{gauthier2024aider} takes a git-native approach:
\begin{itemize}
    \item Automatic commits for each change
    \item Repository map for efficient context selection
    \item Architect/editor model separation for planning vs. implementation
    \item Support for multiple LLM providers
\end{itemize}

\subsubsection{Other Notable Systems}

\textbf{AutoCodeRover}~\citep{zhang2024autocoderover} uses program structure analysis (AST) rather than text search to locate relevant code, improving precision on bug localization.

\textbf{Agentless}~\citep{xia2024agentless} challenges the complexity of agent architectures, showing that a simple two-phase approach (localization then repair) can achieve competitive results with lower cost.

\subsection{Comparison}

\begin{table}[h]
\centering
\caption{Comparison of autonomous coding agents}
\label{tab:agents}
\begin{tabular}{lccccc}
\toprule
\textbf{System} & \textbf{Open} & \textbf{Model} & \textbf{Interface} & \textbf{SWE-bench} \\
\midrule
SWE-agent & Yes & Various & CLI & 12--18\% \\
Devin & No & Proprietary & Web & Claimed high \\
OpenHands & Yes & Various & Web/CLI & 15--25\% \\
Claude Code & No & Claude & CLI & Not published \\
Aider & Yes & Various & CLI & Benchmark varies \\
AutoCodeRover & Yes & Various & CLI & 19--22\% \\
Agentless & Yes & Various & CLI & 18--27\% \\
\bottomrule
\end{tabular}
\end{table}

\subsection{Common Challenges}

Autonomous agents face several persistent challenges:

\textbf{Long-horizon reliability}: Success rates decrease dramatically on tasks requiring many steps. Error cascades are common.

\textbf{Error recovery}: Agents often struggle to recover from mistakes, sometimes entering loops of repeated failed attempts.

\textbf{Security}: Agents with shell access pose security risks. Sandboxing and permission systems are essential but imperfect.

\textbf{Cost}: Agent runs can be expensive, with complex tasks requiring many LLM calls.
